\section{Knowledge Graph Extraction v2}
\label{sec:kg_extraction}

Beyond flat memory storage, MemSt incorporates \textbf{Knowledge Graph Extraction v2} (KGv2), a structured knowledge representation system that extracts entities and relationships from conversations using LLM-powered extraction.

\subsection{Design Motivation}

While tiered memory excels at storing and retrieving factual statements, it lacks explicit relational structure. Consider the following conversation:

\begin{quote}
\textit{``Alice moved to San Francisco last year. She works at OpenAI as a research scientist. Her manager is Bob, who previously worked at Google.''}
\end{quote}

A flat memory system might store three separate facts:
\begin{itemize}
    \item Alice moved to San Francisco in 2022
    \item Alice works at OpenAI
    \item Bob is Alice's manager
\end{itemize}

But the \textit{relationships} between these facts---that Bob works at OpenAI, that Alice's move and job change are related, that Bob's Google experience influences his management---are implicit at best.

KGv2 makes these relationships explicit, enabling:
\begin{enumerate}
    \item \textbf{Multi-hop reasoning}: Answering ``Where did Alice's manager work before?'' requires traversing Alice $\rightarrow$ manager $\rightarrow$ Bob $\rightarrow$ previously worked at $\rightarrow$ Google.
    \item \textbf{Temporal reasoning}: Tracking when relationships were established and how they evolved.
    \item \textbf{Conflict detection}: Identifying when new information contradicts existing relationships.
\end{enumerate}

\subsection{Architecture Overview}

KGv2 follows a modular pipeline architecture:

\begin{enumerate}
    \item \textbf{Extraction}: LLM extracts entities and relationship triplets from text
    \item \textbf{Linking}: Mentions of the same entity are resolved to canonical nodes
    \item \textbf{Storage}: Entities and relationships are persisted to the backend
    \item \textbf{Retrieval}: Query-time subgraph extraction for context assembly
\end{enumerate}

\subsection{Multi-Backend Storage}

A key innovation of KGv2 is pluggable storage backends:

\subsubsection{SQLite Backend}

The default backend uses SQLite with the schema:

\begin{verbatim}
entities(id, name, entity_type, embedding, created_at)
relationships(id, source_id, target_id, relation_type, 
              confidence, created_at, valid_until)
ontologies(id, top_category, first_category, 
           second_category, description)
\end{verbatim}

SQLite is lightweight, requires no external services, and is suitable for development and edge deployment.

\subsubsection{DuckDB Backend}

For production analytics workloads, KGv2 supports DuckDB:

\begin{itemize}
    \item Vectorized query execution for analytical workloads
    \item Efficient aggregations and graph analytics (PageRank, community detection)
    \item Parquet export for integration with data lakes
\end{itemize}

The backend is selected at compile time via Cargo features:

\begin{verbatim}
# SQLite (default)
memst-extract-v2 = { path = "../memst-extract-v2" }

# DuckDB (production)
memst-extract-v2 = { path = "../memst-extract-v2", 
                     default-features = false, 
                     features = ["duckdb"] }
\end{verbatim}

\subsection{Ontology System}

KGv2 includes a comprehensive ontology system supporting 80+ intelligence domains:

\begin{table}[ht]
\centering
\caption{Sample Ontology Categories}
\begin{tabular}{@{}lll@{}}
\toprule
\textbf{Top Category} & \textbf{First Category} & \textbf{Second Category} \\
\midrule
Domain Intelligence & Tech Intelligence & Artificial Intelligence \\
Domain Intelligence & Tech Intelligence & Semiconductor \\
Domain Intelligence & Business Intelligence & Market Analysis \\
Domain Intelligence & Business Intelligence & Competitive Intelligence \\
Social Intelligence & Social Media & Trend Analysis \\
Social Intelligence & Public Opinion & Sentiment Monitoring \\
Geopolitical & Regional & Asia-Pacific \\
Geopolitical & Regional & European Union \\
\bottomrule
\end{tabular}
\label{tab:ontologies}
\end{table}

Each ontology defines:
\begin{itemize}
    \item Expected entity types (Person, Organization, Location, Event, Technology)
    \item Relationship vocabulary (works\_at, located\_in, acquired, partnered\_with)
    \item Extraction prompts tuned for the domain
\end{itemize}

Users can load custom ontologies via JSON schema:

\begin{verbatim}
{
  "top_category": "Custom Domain",
  "first_category": "Research",
  "second_category": "Machine Learning",
  "chinese_name": "Research-ML",
  "english_name": "Research-ML",
  "overview": "ML research tracking"
}
\end{verbatim}

\subsection{Entity Extraction Pipeline}

The extraction pipeline uses a two-stage LLM approach:

\textbf{Stage 1: Entity Detection}
\begin{enumerate}
    \item Input: Conversation text + ontology context
    \item LLM identifies entity mentions with types and confidence scores
    \item Output: List of $(\text{mention}, \text{type}, \text{confidence})$ tuples
\end{enumerate}

\textbf{Stage 2: Relationship Extraction}
\begin{enumerate}
    \item Input: Detected entities + full text
    \item LLM identifies relationships as $(\text{source}, \text{relation}, \text{target})$ triplets
    \item Confidence scoring based on explicit vs. implicit mentions
\end{enumerate}

The extraction is performed incrementally: new messages are processed and merged into the existing knowledge graph rather than regenerating from scratch.

\subsection{Entity Linking and Canonicalization}

A critical challenge is resolving when two mentions refer to the same entity:

\begin{quote}
\textit{``OpenAI released GPT-4. The San Francisco-based company also announced...''}
\end{quote}

Here ``OpenAI'' and ``The San Francisco-based company'' refer to the same organization. KGv2 implements:

\begin{itemize}
    \item \textbf{Exact matching}: Name string normalization and comparison
    \item \textbf{Embedding similarity}: Cosine similarity of entity embeddings
    \item \textbf{Context clustering}: Entities appearing in similar contexts are candidates for merging
    \item \textbf{LLM verification}: For high-stakes merges, an LLM confirms the equivalence
\end{itemize}

Linked entities are stored with a canonical ID, while surface forms are preserved for provenance.

\subsection{Conflict Detection and Resolution}

When new information contradicts existing relationships, KGv2 applies temporal logic:

\begin{itemize}
    \item New relationships with later timestamps supersede older ones
    \item Contradicted relationships are marked $\text{valid\_until} = t_{\text{new}}$ rather than deleted
    \item This enables temporal queries: ``Where did Alice work in 2021?''
\end{itemize}

For complex conflicts (e.g., contradictory information from different sources), the system can escalate to an LLM for resolution or flag for human review.

\subsection{Integration with Tiered Memory}

KGv2 operates in concert with the four-tier memory system:

\begin{itemize}
    \item Working memory contains the ``hot'' subgraph most relevant to current context
    \item Short-term memory maintains recent entity observations
    \item Long-term memory stores the canonical entity catalog
    \item Archival memory preserves historical relationship states
\end{itemize}

During context assembly, MemSt can include both flat memories and structured subgraphs, with the LLM determining how to combine them for answering the query.
